\slajdid{08-s058}
\begin{frame}[label=08-s058]{Standardna situacija}
\gusttekst\setlength{\parskip}{2pt}\setlength{\abovedisplayskip}{3pt}\setlength{\belowdisplayskip}{3pt}
\begin{columns}[c,onlytextwidth]\column{.44\textwidth}\centering\input{slike/tikz/s058_rc.tex}
\column{.54\textwidth}\centering\begin{tikzpicture}[x=.78cm,y=.6cm,font=\sitneoznake,>=Latex,remember picture]
\draw[->](-1.5,0)--(4.8,0)node[below]{$t$};\draw[->](0,-.1)--(0,2.4)node[left]{$u_i(t)$};
\draw(-1,.3)--(.8,.3)--(.8,1.9)--(3.3,1.9)--(3.3,.3)--(4.4,.3);\draw[dotted](3.3,0)--(3.3,.3);\node[below]at(3.3,0){$t_2$};\draw[dotted](.8,0)--(.8,.3);\node[below]at(.8,0){$t_1$};\node[above]at(2.3,1.9){$=V_{OH}$};\node[above]at(-.5,.3){$=V_{OL}$};\coordinate(s58end)at(4.4,.3);\end{tikzpicture}
\end{columns}
\[u_o(t)=u_o(\infty)+\bigl(u_o(t_1^+)-u_o(\infty)\bigr)e^{-\frac{t-t_1}{\tau}}\quad za\ t\geq t_1\]
\begin{columns}[T,onlytextwidth]\column{.44\textwidth}
1. $\tau$\quad $\tau=C*R$\par
2. $u_o(t_1^+)$\par
Kako je izlazni napon istovremeno i napon na kondenzatoru važi
\[u_0(t_1^+)=u_c(t_1^+)=u_c(t_1^-)\]
to je\qquad $u_0(t_1^+)=V_{OL}$
\column{.54\textwidth}
3. $u_o(\infty)$\par
Kada se završe svi prelazni procesi $i_c(\infty)=0=i_r(\infty)$. Kako je u beskonačnosti napon na ulazu \tikz[remember picture,baseline=(s58text.base)]{\node[inner sep=0pt](s58text){$V_{OL}$};} i struja kroz otpornik jednaka nuli\tikz[remember picture,overlay]{\draw[->,DEplava](s58text.south)--++(0,-3mm)-|([xshift=21mm,yshift=-3mm]s58end)--([xshift=21mm]s58end)--(s58end);}
\[u_0(\infty)=u_i(\infty)-Ri_R(\infty)=V_{OL}\]
\end{columns}
\[u_o(t)=u_o(\infty)+\bigl(u_o(t_0^+)-u_o(\infty)\bigr)e^{-\frac{t-t_1}{\tau}}\quad za\ t\geq t_1\]
\[\alert{u_o(t)=V_{OL}+(V_{OL}-V_{OL})e^{-\frac{t-t_1}{\tau}}=V_{OL}\quad za\ t\geq t_1}\]
\hfill ????????????
\end{frame}
